\documentclass[sigconf, screen, review, nonacm, anonymous]{acmart}
%\documentclass[sigconf]{acmart}

% Hashimoto added
\sloppy % 無理に詰めてoverfullになるのを防ぐ。
\usepackage{xcolor}
\newcommand{\TODO}[1]{\textcolor{red}{\textbf{TODO: #1}}}
\usepackage{booktabs}
\usepackage{comment}

\usepackage{multirow} 


%%
%% end of the preamble, start of the body of the document source.
\begin{document}
%\balance
%%
%% The "title" command has an optional parameter,
%% allowing the author to define a "short title" to be used in page headers.
\title{Manga109Script: Manga Script Dataset for Script2Layout -supplementary material-}

\maketitle
\appendix
\renewcommand\thefigure{\thesection.\arabic{figure}}    

\section{Training Conditions}
Model training was performed using the \texttt{transformers} and \texttt{peft} libraries with the following configuration:

\begin{itemize}
    \item \textbf{Dataset:}
    \begin{itemize}
        \item The \texttt{Manga109ScriptDataset} is used with:
        \begin{itemize}
            \item Maximum token length: 4096,
            \item The data sequence is divided into 32 parts.
        \end{itemize}
    \end{itemize}
    \item \textbf{LoRA Configuration:}
    \begin{itemize}
        \item Reduction factor $r = 8$,
        \item \texttt{lora\_alpha} = 16,
        \item Target modules: \texttt{q\_proj} and \texttt{v\_proj},
        \item LoRA dropout: 0.1,
        \item Bias setting: \texttt{none},
        \item Task type: \texttt{CAUSAL\_LM}.
    \end{itemize}
    \item \textbf{Training Parameters:}
    \begin{itemize}
        \item \textbf{Batch Size:} 1 per device for both training and evaluation,
        \item \textbf{Gradient Accumulation:} 32 steps,
        \item \textbf{Epochs:} 20,
        \item \textbf{Optimizer:} \texttt{adamw\_torch\_fused},
        \item \textbf{Learning Rate:} $2 \times 10^{-4}$,
        \item \textbf{LR Scheduler:} Cosine scheduler,
        \item \textbf{FP16:} Enabled,
        \item \textbf{Logging:} Every 100 steps,
        \item \textbf{Evaluation Strategy:} Evaluation performed every 100 steps with data grouped by length,
    \end{itemize}
\end{itemize}

The training is managed through the \texttt{Trainer} class provided by the \texttt{transformers} library.\footnote{\url{https://pypi.org/project/transformers/}~(cited 2025-04-11)} Finally, after training, the model is saved along with the LoRA adapter parameters (saved separately) and the tokenizer.

\section{Additional Qualitative Results}
In addition to Fig. 3 of the main paper, we introduce a failure case of script2layout models in Fig. \ref{fig:quantitative_results}.
An extreme case is the excessive enlargement of speech balloons. This occurs when there is a large amount of dialogue, and the Script2Layout model directly takes the dialogue content as input.


\begin{figure*}[h]
  \includegraphics[width=\textwidth]{img/failure.pdf}
  \caption{
Failure cases of generated layouts. Each row shows one example: the real manga layout (left), followed by predictions from existing methods and Script2Layout models as in Fig. 3 of the main paper (Source: *Touta Mairimasu* by Shinji Saijo)
}
  \Description{}
  \label{fig:quantitative_results}
\end{figure*}

\end{document}